%% STEP 5 — 원고 핵심 섹션 재작성
%% 파일: manuscript_v2_abstract_intro.tex
%% 적용 방법: 아래 각 섹션을 paper/main_final.tex에서 해당 섹션으로 교체할 것.
%%
%% 재작성 원칙:
%% - Abstract 첫 3문장 안에 novelty (failure regime + insight) 가시화
%% - Introduction: 아키텍처 설명 전에 "왜 새로운가"를 먼저 제시
%% - Contribution: problem/method/evidence 구조로 3개
%% - Related Work: 나열 → 각 문단 끝에 "따라서 ___가 미해결" 명시

%%==============================================================================
%% [1] ABSTRACT — 완전 재작성
%%==============================================================================

%% 구조: 문제→기존한계→핵심아이디어→근거1→근거2→함의 (6문장)

\begin{abstract}
RAG-based expert systems route queries by domain identity alone, ignoring the
reasoning function a query demands; we show this single-axis design fails under
two structural conditions---functional mismatch and cross-domain evidence
contamination---neither of which can be resolved by improving routing accuracy alone.
We propose \textbf{NexRoute}, a dual-axis framework whose central finding is that
\emph{structural wiring}---binding each routing decision jointly to a retrieval
partition and a function-conditioned prompt scaffold---matters more than routing
probability optimization: oracle-domain routing (the accuracy upper bound) degrades
generation quality by $-1.72$~pp ROUGE-L relative to predicted routing ($p{<}0.001$),
while LoRA-library methods with comparable routing entropy lag NexRoute by
${\geq}18$~pp ROUGE-L on cross-domain queries despite routing to the same adapter pool.
A jointly trained DeBERTa-large encoder produces a domain score and a function label
in a single forward pass, directing each query to a domain-exclusive retrieval
partition, a LoRA-specialized adapter, and a reasoning-aligned prompt scaffold
simultaneously; cross-domain queries additionally activate a dedicated interdisciplinary
adapter with evidence drawn separately from two domain partitions.
On a 300-sample benchmark fully independent of LoRA adapter training, this structural
alignment enables NexRoute 7B to exceed all 7B baselines by ${\geq}6.06$~pp
CROSS-domain ROUGE-L (Holm-corrected $p{<}0.001$) and attain the lowest CROSS-domain
sentence-level NLI contradiction rate among 7B models ($3.6\%$), while
operating within 14.8~GB of VRAM.
Ablation confirms a clear division of labour: removing retrieval partitioning
elevates CROSS-domain NLI contradiction rate by $+3.20$~pp ($p{=}0.045$), while
removing function scaffolding costs $1.32$~ROUGE-L ($p{<}0.001$), with each axis
governing a structurally distinct quality dimension.
These results reframe the design objective for cross-domain expert systems: rather
than asking which expert to activate, the productive question is which evidence
scope and reasoning format to bind to each query---and these must be determined jointly.
\end{abstract}

%%==============================================================================
%% [2] INTRODUCTION — 완전 재작성
%%==============================================================================

\section{Introduction}
\label{sec:introduction}

%%--- 단락 1: 문제 설정 (failure regime 먼저)
Large Language Models (LLMs) increasingly serve professional domains where queries
require not only domain-specific knowledge but a specific reasoning operation:
a physician's question may demand causal explanation, while a lawyer's may demand
definitional recall over the same underlying concepts
\citep{brown2020gpt,singhal2023large,guha2023legalbench}.
Existing RAG-based expert systems address the first requirement---domain
specialization---through LoRA adapters and retrieval augmentation
\citep{hu2022lora,dou2024loramoe}, but structurally ignore the second.
We identify this omission as the source of two distinct, coexisting failure modes.

%%--- 단락 2: Failure Mode 1 (functional mismatch)
The first failure mode is \textbf{functional mismatch}: a domain-correct expert
nonetheless fails because its generation prompt provides no structural signal
for the reasoning operation the query demands---factual recall, procedural
explanation, causal justification, or comparison.
A LoRA adapter optimized for causal legal reasoning produces suboptimal responses
when activated for a definitional query in the same domain, because the scaffold
is silent on reasoning modality \citep{singhal2023large,guha2023legalbench}.
Single-axis domain routing is structurally blind to this functional heterogeneity.

%%--- 단락 3: Failure Mode 2 (cross-domain contamination)
The second failure mode is \textbf{cross-domain evidence contamination}: when
interdisciplinary queries retrieve from a unified corpus, lexical cues shared
across professional domains---``liability'' appearing in both medical consent
records and legal statutes; ``mechanism'' in both pharmacology and engineering
procedure---cause retrieval to return domain-incongruent passages.
Our ablation quantifies this: removing domain-partitioned retrieval from NexRoute
specifically elevates CROSS-domain NLI contradiction rate by $+3.20$~pp
($3.50\%{\to}6.70\%$, $p{=}0.045$), while single-domain contradiction rates
are unchanged---isolating retrieval-scope control as the operative mechanism,
not a general quality improvement (Figure~\ref{fig:motivation}).

%%--- 단락 4: 핵심 통찰 (counter-intuitive finding, first)
A central and counter-intuitive finding motivates our approach:
optimizing routing accuracy toward its oracle upper bound \emph{does not improve}---and
in fact degrades---generation quality.
Oracle-domain routing (gold domain labels, accuracy upper bound) reduces ROUGE-L
by $-1.72$~pp overall and $-1.96$~pp single-domain relative to NexRoute predicted routing
($p{<}0.001$ for both; Table~\ref{tab:domainrouter}).
Furthermore, LoRA-library routing methods (Arrow Routing, PHATGOOSE) achieve routing
entropy comparable to NexRoute on cross-domain queries ($91.5\%$, $100.0\%$ vs.\
$92.5\%$ of $\mathcal{H}_{\max}$), yet lag NexRoute by ${\geq}18$~pp ROUGE-L on the same queries.
Together, these results establish that \textbf{routing probability optimization is not
the productive design choice}: what routing selects---the evidence scope and reasoning
format---determines quality, not how confidently it selects.

%%--- 단락 5: NexRoute 제안 (architecture AFTER insight)
To instantiate this design principle, we propose \textbf{NexRoute}, a dual-axis
framework that simultaneously determines (1) which domain-exclusive retrieval
partition and LoRA adapter to apply, and (2) which function-conditioned prompt
scaffold to use.
A shared DeBERTa-large encoder produces both decisions in one forward pass:
a \textbf{domain axis} ($R_D$: \texttt{LAW}, \texttt{MED}, \texttt{SOC},
\texttt{STEM}, \texttt{CROSS}) binds the query to a retrieval partition and expert
adapter, while a \textbf{function axis} ($R_F$: \texttt{What}, \texttt{How},
\texttt{Why}, \texttt{Compare}) binds the query to a reasoning-aligned prompt scaffold.
The intersection of these decisions---the \textit{nexus point}---synchronizes evidence
scope and generation format, addressing both failure modes in one architectural step.
Cross-domain queries activate a dedicated \texttt{CROSS} adapter and retrieve
evidence independently from two domain partitions, restricting evidence mixing by
construction ($\mathrm{IDI}(q){=}0$, Eq.~\ref{eq:partition}).

%%--- 단락 6: 기여 3개 (problem/method/evidence 구조)
Three contributions follow:

\begin{itemize}
    \item \textbf{(Problem)} We formally define two failure regimes in single-axis
    routing---functional mismatch and cross-domain contamination---and quantify their
    structural effects: removing function scaffolding costs $1.32$~ROUGE-L
    ($p{<}0.001$) and removing retrieval partitioning elevates CROSS-domain NLI
    contradiction rate by $+3.20$~pp ($p{=}0.045$), with each component governing
    a structurally distinct quality dimension that the other cannot compensate.

    \item \textbf{(Method)} We introduce the \emph{nexus point}---a jointly trained
    dual-axis decision that simultaneously binds domain routing to retrieval partition
    and adapter selection, and function routing to prompt scaffold selection, in a
    single synchronized forward pass.
    Cross-domain detection requires no dedicated fifth output class: it is inferred
    from the joint activation pattern of top-two domain scores via a threshold
    criterion (Eq.~\ref{eq:cross_detect}), enabling soft interdisciplinary query
    detection.
    Theoretical analysis (Section~\ref{sec:theory}) establishes that single-axis
    routing has strictly lower expected quality under function heterogeneity, a bound
    confirmed by our controlled ablation.

    \item \textbf{(Evidence)} We demonstrate through three converging lines of
    evidence that structural wiring---not routing optimization---explains quality
    gains: (i)~oracle-domain routing degrades ROUGE-L by $-1.72$~pp ($p{<}0.001$);
    (ii)~Arrow Routing and PHATGOOSE, despite routing entropy comparable to NexRoute,
    lag by ${\geq}18$~pp ROUGE-L; and (iii)~ablation reveals complementary,
    non-overlapping quality contributions of each structural component.
    On the contamination-free 300-sample benchmark, structural alignment alone
    enables NexRoute 7B to exceed all 7B baselines by ${\geq}6.06$~pp CROSS-domain
    ROUGE-L (Holm-corrected $p{<}0.001$, 87.4\% routing accuracy on 500-sample test).
\end{itemize}

%%==============================================================================
%% [3] RELATED WORK — 재작성 방향 (기존 텍스트 보강)
%%==============================================================================
%% 각 subsection 끝에 "따라서 ___가 미해결" 문장 추가
%% 나열 대신 비교-대조 서술로 전환

%% 3.1 Mixture-of-Experts — 추가할 마지막 문장:
%%
%%  NexRoute departs from this line by operating at the query level (not the token level)
%%  and binding routing decisions to retrieval scope and prompt format simultaneously.
%%  The question these MoE systems leave open is not \emph{which} expert to activate,
%%  but \emph{what} the activation should control downstream---a gap NexRoute addresses.

%% 3.2 LoRA-MoE — 추가할 마지막 문장:
%%
%%  The critical gap these systems share is the decoupling of adapter selection
%%  from retrieval scope and prompt structure: routing determines the expert,
%%  but not the evidence the expert receives or the reasoning format it operates under.
%%  NexRoute closes this gap by wiring adapter selection to domain-exclusive retrieval
%%  and function classification to scaffold selection, in one synchronized step.

%% 3.3 RAG and Multi-Domain — 추가할 마지막 문장:
%%
%%  Thus, the retrieval bottleneck for cross-domain queries---evidence contamination
%%  from lexically similar but domain-incongruent passages---remains open in all
%%  standard RAG variants, as retrieval operates over a unified index without
%%  domain-exclusive partitioning.

%% 3.4 Query Routing — 추가할 마지막 문장 (기존 "Structural delta" 단락 강화):
%%
%%  Our oracle routing experiment directly addresses this shared gap: it shows
%%  that routing accuracy optimization is insufficient and potentially harmful,
%%  because quality arises from the structural resources routing selects,
%%  not from the routing decision's confidence.
%%  No prior system has demonstrated this empirically; NexRoute provides the
%%  first controlled evidence that structural wiring supersedes routing optimization.

%%==============================================================================
%% [4] EXPERIMENTS — 재배치 지침
%%==============================================================================
%% 현재: 실험 전체가 "성능 비교표" 위주
%% 수정: E1 (가설 검증)을 Section 4.2로 이동, 기존 routing accuracy가 Section 4.1이 됨
%%
%% 새로운 Experiments 구조:
%% 4.1 Experimental Setup
%% 4.2 Hypothesis Test: Structural Wiring over Routing Accuracy [NEW: E1]
%%     → 이 섹션을 논문에서 가장 먼저 배치 (성능표 앞)
%%     → Oracle routing < Predicted routing: 핵심 발견 여기서 먼저 소개
%% 4.3 Routing System Accuracy (기존 Section 4.2)
%% 4.4 Comprehensive Generation Quality (기존 Section 4.3)
%% 4.5 Factuality Assessment (기존 Section 4.4)
%% 4.6 LAW Adapter Correction (기존 Section 4.5)
%% 4.7 Domain Router Necessity [통합: 기존 Section 4.6]
%% 4.8 External Benchmark Evaluation (기존 Section 4.7)
%%
%% 이 재배치 이유: "성능이 좋다" 대신 "왜 새로운가"를 먼저 보여주는 구조

%%==============================================================================
%% [5] DISCUSSION — 추가할 Limitation 3개 (간결 버전)
%%==============================================================================
%%
%% 심사 통과율을 올리는 limitation 서술 (정직 + 구체적):
%%
%% L1. Empirical scope: 모든 결과는 MMLU-derived corpora와 Qwen2.5 backbone으로 제한됨.
%%     Cross-backbone validation (LLaMA-3 8B, Mistral 7B)과 human-authored 참조
%%     답변이 없으면 broader generalizability를 주장할 수 없음.
%%
%% L2. Theoretical gap: Proposition 1은 scaffold benefit의 existence만 보이며,
%%     magnitude는 dataset-specific임. Oracle routing < predicted routing의 메커니즘은
%%     3가지 가설(adapter distribution match, retrieval composition, query phrasing)로
%%     제시되나 controlled verification이 부재함.
%%
%% L3. Retrieval quality: IDI(q)=0이 정의상 보장되지만 retrieved passages의
%%     topical relevance (Recall@K, nDCG)는 미평가. Domain membership ≠ topical relevance.

%%==============================================================================
%% [6] Highlights (revision용)
%%==============================================================================

%% \begin{highlights}
%%     \item Oracle-domain routing (accuracy upper bound) degrades ROUGE-L by 1.72pp,
%%           demonstrating that structural wiring supersedes routing accuracy optimization.
%%     \item Removing domain-partitioned retrieval specifically elevates CROSS-domain
%%           NLI contradiction rate (+3.20pp, p=0.045), isolating retrieval scope control
%%           as the mechanism addressing cross-domain contamination.
%%     \item Arrow and PHATGOOSE achieve comparable routing entropy to NexRoute yet lag
%%           by ≥18pp ROUGE-L on CROSS queries, confirming structural design—not
%%           routing probability—as the differentiating factor.
%%     \item NexRoute 7B exceeds all 7B baselines by ≥6.06pp CROSS-domain ROUGE-L
%%           (p<0.001) on a contamination-free 300-sample benchmark within 14.8GB VRAM.
%% \end{highlights}
